% Chapter 3, section 10: Retries and dead letters
% Write in plain LaTeX. Code of a guide step: \guide{<step>}; commands: \begin{shell}...\end{shell};
% a screenshot: \screenshot{file}{caption}; a gap in the draft: \placeholder{...}.
% See book/latex/queen.sty for the full list.
%
% Draft by Claude, 2026-10-24. The code is cut from examples/tutorials/js/book-use-cases.mjs, which
% checks every result quoted here against a 2.6.0 node.

\section{Retries and dead letters}

When a booking is confirmed, send \texttt{booking.confirmed} to the customer's URL. Their server
fails sometimes, and a crash at the wrong moment can make you call it twice.

One message per call, in a queue with retries (\texttt{retryLimit: 5} and a \texttt{retryDelay}).
The handler throws when the call fails: the message comes back after the delay, and goes to the
dead-letter queue after the last try. The message's \texttt{transactionId} travels as the
\texttt{Idempotency-Key}, so the receiver can drop a second call.

\queenfig{webhook-retries}{A failed call comes back after the delay; the key makes it count once.}

\guide{uc-webhook}

The receiver failed the first call. The message came back after the delay, the second call worked,
and the receiver kept the event once, under \texttt{evt-9001}.

\why{the call happens outside any commit, so it can happen twice: a crash after the call and before
the ack sends it again. The receiver dropping the second call by its key is what makes it count
once. Queen gives the stable key, the retries and the dead-letter queue.}

One thing to know: a failed message holds its partition for the delay. Use one partition per
customer, so a customer whose server is down waits on its own retries alone.
